One polynomial satisfying the three conditions is $P_0(x)=x^2-9x+20$. For any other such polynomial, $P-P_0$ vanishes at 2 and 6, so
\[P(x)=P_0(x)+(x-2)(x-6)(ux+v).\]
The condition $P(1)=P(8)$ gives $5(u+v)=12(8u+v)$, hence $v=-13u$. Thus all degree-three solutions are
\[P(x)=x^2-9x+20+u(x-2)(x-6)(x-13),\qquad u\ne0.\]
For $P(a)$ to be independent of $u$, $a$ must be 2, 6, or 13. The first two are excluded, so $a=13$ and $b=P_0(13)=72$. Therefore $a+b=\boxed{85}$.
